\documentclass[11pt]{article}

\usepackage[margin=1in]{geometry}
\usepackage{amsmath,amssymb,amsfonts}
\usepackage{booktabs}
\usepackage{longtable}
\usepackage{array}
\usepackage{enumitem}

\title{Geometric \& Topological Methods in Atmospheric Dynamics\\
\large Reference Card -- UAH Graduate Course}
\date{}

\begin{document}

\maketitle

%===========================================================
\section*{I. Differential Geometry}

\begin{longtable}{>{\bfseries}p{0.28\textwidth}p{0.67\textwidth}}
Object & Definition\\
\midrule
Manifold $M$ &
Smooth space locally homeomorphic ($\cong$) to $\mathbb{R}^n$.\\

Tangent space $T_pM$ &
Vector space of velocity vectors (derivations) at $p\in M$.\\

Cotangent space $T_p^*M$ &
Dual space to $T_pM$; contains $1$-forms (covectors).\\

$k$-form &
Smooth, totally antisymmetric multilinear functional acting on $k$ tangent vectors.
\end{longtable}

\subsection*{Graded Algebra of Forms ($\Omega^*(M)$)}

\begin{itemize}[leftmargin=1.5cm]
\item \textbf{0-form}: scalar field ($p$, $T$, $\theta$)
\item \textbf{1-form}: $v$ --- velocity/circulation degrees of freedom
\item \textbf{2-form}: $\omega=dv$ --- vorticity/flux degrees of freedom
\end{itemize}

\subsection*{Key Operators}

\begin{longtable}{>{\bfseries}p{0.22\textwidth}p{0.28\textwidth}p{0.42\textwidth}}
Operator & Action & Key Property\\
\midrule

Exterior derivative $d$
&
$\Omega^k(M)\rightarrow\Omega^{k+1}(M)$
&
$d^2=0$ (nilpotent; boundaries have no boundary).\\

Wedge product $\wedge$
&
$\Omega^p(M)\times\Omega^q(M)\rightarrow\Omega^{p+q}(M)$
&
Graded commutative:
\[
\alpha\wedge\beta=(-1)^{pq}\beta\wedge\alpha.
\]
\\

Hodge star $\star$
&
$\Omega^k(M)\rightarrow\Omega^{n-k}(M)$
&
Isomorphism; explicitly metric-dependent.\\

Interior product $\iota_u$
&
$\Omega^k(M)\rightarrow\Omega^{k-1}(M)$
&
Contraction with vector field $u$.\\

Lie derivative $\mathcal{L}_u$
&
$\Omega^k(M)\rightarrow\Omega^k(M)$
&
Measures transport/deformation along the flow.
\end{longtable}

\paragraph{Cartan's Magic Formula}

\[
\boxed{
\mathcal{L}_u\alpha
=
d(\iota_u\alpha)
+
\iota_u(d\alpha)
}
\]

\paragraph{Stokes' Theorem}

\[
\boxed{
\int_{\partial M}\omega
=
\int_M d\omega
}
\]

%===========================================================
\section*{II. Geometric Fluid Dynamics (Incompressible Euler)}

\subsection*{Vector-Invariant Momentum Equation}

\[
\boxed{
\partial_t v
+
\iota_u\omega
+
dB
=
0\,
\qquad
\omega=dv
}
\]

\paragraph{Bernoulli Function}

\[
B
=
\frac12\iota_u v
+
p
+
\Phi .
\]

\noindent
($\tfrac12|u|^2$ is written coordinate-free as
$\tfrac12\iota_u v$.)

\paragraph{Mass Conservation}

\[
d(\star v)=0.
\]

\paragraph{Vorticity Transport}

\[
\partial_t\omega+\mathcal L_u\omega=0
\]

or equivalently,

\[
\partial_t\omega+d(\iota_u\omega)=0.
\]

\subsection*{Conserved Invariants}

\begin{longtable}{>{\bfseries}p{0.22\textwidth}p{0.28\textwidth}p{0.42\textwidth}}
Invariant & Expression & Physical Mechanism\\
\midrule

Energy
&
\[
E=\frac12\int_M v\wedge\star v
\]
&
The Lamb vector satisfies
\[
\iota_u(\iota_u\omega)=0\,
\]
so it performs no work.\\

Kelvin circulation
&
\[
\frac{d}{dt}\oint_{c(t)}v=0
\]
&
Direct consequence of Stokes' theorem and
\[
\mathcal L_u v=-dB.
\]
\\

Helicity (3D)
&
\[
H=\int_M v\wedge dv
\]
&
Topological measure of vortex-line linking and knotting.\\

Potential vorticity
&
\[
q=\frac{\star(dv\wedge d\theta)}{\rho}
\]
&
Material invariant
\[
\frac{Dq}{Dt}=0
\]
for stratified adiabatic flow.\\

Enstrophy (2D)
&
\[
Z=\frac12\int_M\omega\wedge\star\omega
\]
&
Conserved in two dimensions; drives the inverse energy cascade.
\end{longtable}

%===========================================================
\section*{III. Discrete Exterior Calculus (DEC)}

\paragraph{Philosophy}

Separate purely combinatorial topology (incidence matrices) from
geometric measurements (lengths, areas, volumes).

\begin{longtable}{>{\bfseries}p{0.28\textwidth}p{0.65\textwidth}}
Element & Description\\
\midrule

Primal complex &
Simplicial complex (e.g.\ Delaunay triangulation); carries primal
$k$-cochains.\\

Dual complex &
Circumcentric dual (e.g.\ Voronoi tessellation); carries dual
$(n-k)$-cochains.\\

Coboundary operator $D$ &
Discrete analogue of $d$ defined entirely from incidence matrices.
Exactly satisfies
\[
D^2=0.
\]
\\

Discrete Hodge star $M_k$ &
Diagonal geometric mass matrix mapping primal to dual cochains;
the only source of numerical error.\\

Whitney forms &
Interpolate discrete cochains back to continuous differential forms.
\end{longtable}

\paragraph{Discrete Energy Identity}

\[
\boxed{
\langle v\,\iota_u(\widetilde Dv)\rangle=0
}
\]

Preserving skew-symmetry of the convective operator guarantees
unconditional structural stability and eliminates artificial numerical
vorticity dissipation.

\paragraph{Atmospheric Applications}

MPAS, ICON, mimetic finite differences, and compatible geometric finite
elements.

%===========================================================
\section*{IV. Dynamical Systems \& Bifurcation Theory}

\subsection*{Local Bifurcations}

\begin{longtable}{>{\bfseries}p{0.18\textwidth}p{0.28\textwidth}p{0.46\textwidth}}
Type & Normal Form & Atmospheric Example\\
\midrule

Saddle-node &
\[
\dot x=\mu-x^2
\]
&
Climate tipping points; fold catastrophes.\\

Pitchfork &
\[
\dot x=\mu x-x^3
\]
&
Symmetry breaking (Hadley circulation, monsoon asymmetry).\\

Transcritical &
\[
\dot x=\mu x-x^2
\]
&
Threshold crossing to instability.\\

Hopf &
\[
\dot z=(\mu+i\omega)z-z|z|^2
\]
&
Onset of oscillatory atmospheric behavior.
\end{longtable}

\paragraph{Global Bifurcations}

Homoclinic and heteroclinic connections, period-doubling cascades,
torus breakdown, and interior/boundary crises.

\subsection*{Catastrophe Theory}

\begin{longtable}{>{\bfseries}p{0.18\textwidth}p{0.12\textwidth}p{0.18\textwidth}p{0.40\textwidth}}
Catastrophe & Codim. & Controls & Atmospheric Relevance\\
\midrule

Fold & 1 & $\mu$ &
Abrupt climate or ecosystem transitions.\\

Cusp & 2 &
$(\alpha,\beta)$ &
Bistability and hysteresis; stratospheric warmings and SBL transitions.\\

Swallowtail & 3 &
Three parameters &
Ocean--atmosphere--ice interactions.\\

Butterfly & 4 &
Four parameters &
Structural stability boundaries in highly nonlinear models.
\end{longtable}

%===========================================================
\section*{V. Lagrangian Coherent Structures (LCS)}

\paragraph{Flow Map}

\[
F_{t_0}^{\,t}(x_0)=x(t; t_0, x_0).
\]

\paragraph{Cauchy--Green Tensor}

\[
\boxed{
C=(DF_{t_0}^{\,t})^{T}DF_{t_0}^{\,t}
}
\]

\paragraph{Finite-Time Lyapunov Exponent}

\[
\boxed{
\sigma_{t_0}^{\,t}(x_0)
=
\frac{1}{|t-t_0|}
\ln
\sqrt{\lambda_{\max}(C(x_0))}
}
\]

\subsection*{LCS Taxonomy}

\begin{longtable}{>{\bfseries}p{0.20\textwidth}p{0.22\textwidth}p{0.48\textwidth}}
Type & FTLE Field & Geometry\\
\midrule

Repelling &
Forward-time ridges &
Material lines maximizing normal stretching.\\

Attracting &
Backward-time ridges &
Material lines maximizing normal compression.\\

Elliptic &
Closed nested loops &
Coherent vortex cores.\\

Parabolic &
Shear trenches &
Material transport barriers.
\end{longtable}

\paragraph{Atmospheric Applications}

Atmospheric rivers, hurricane eyewalls, jet-stream filaments,
polar vortex boundaries, and volcanic ash transport.

%===========================================================
\section*{VI. Topological Data Analysis (TDA)}

\paragraph{Persistent Homology Pipeline}

\[
\boxed{
\text{Point Cloud\/Field Data}
\rightarrow
\text{Filtered Complex}
\rightarrow
H_k
\rightarrow
\text{Persistence Diagrams\/Barcodes}
}
\]

\begin{longtable}{>{\bfseries}p{0.20\textwidth}p{0.30\textwidth}p{0.40\textwidth}}
Object & Mathematical Meaning & Atmospheric Meaning\\
\midrule

$\beta_0$ &
Connected components &
Cloud patches or isolated structures.\\

$\beta_1$ &
Independent loops &
Fluid circuits or vortex rings.\\

$\beta_2$ &
Three-dimensional voids &
Enclosed bubbles or cavities.\\

Persistence diagram &
Birth/death point set &
Evolution of geometric structures.\\

Barcode &
Intervals $[b,d]$ &
Feature lifetime. Long bars indicate robust structures.\\

Mapper &
Topological graph &
Reduced state-space representation.
\end{longtable}

\paragraph{Atmospheric Applications}

Convective organization, cloud lifecycle analysis, turbulence
intermittency, structural regime tracking, and radar morphology.

%===========================================================
\section*{VII. Atmospheric Boundary Layer Geometry}

\begin{longtable}{>{\bfseries}p{0.30\textwidth}p{0.62\textwidth}}
Concept & Geometric Interpretation\\
\midrule

Monin--Obukhov similarity &
Invariant coordinate scaling defining a smooth similarity manifold.\\

Stability regimes &
Folded equilibrium surface with structural transitions.\\

Cusp bifurcation &
Explains hysteresis between SBL and VSBL.\\

Richardson number ($Ri$) &
Primary bifurcation parameter on the stability manifold.\\

Slow manifold &
Attracting invariant manifold after gravity-wave adjustment.\\

Geometry-aware closures &
Reynolds stresses represented as geometric invariants tied to metric curvature.
\end{longtable}

\subsection*{Key Dimensionless Ratios}

\[
Ri
=
\frac{
\frac{g}{\theta_v}
\frac{\partial\theta_v}{\partial z}
}{
\left(\frac{\partial u}{\partial z}\right)^2
+
\left(\frac{\partial v}{\partial z}\right)^2
}\,
\qquad
Ro=\frac{U}{fL}\,
\qquad
Fr=\frac{U}{NH}.
\]

%===========================================================
\section*{Suggested Revision Notes for Class Delivery}

\begin{enumerate}
\item
Emphasize that the geometric term
\[
\iota_u\omega
\]
is precisely the coordinate-free analogue of the classical vector identity
\[
u\times(\nabla\times u).
\]
This is often the conceptual bridge for engineering students.

\item
Remind students that the exterior derivative $d$ is completely metric-free,
whereas the Hodge star $\star$ introduces the geometry of the physical
space (flat Cartesian grids, spherical Earth, etc.).
\end{enumerate}

\end{document}
